\documentclass[12pt,a4paper]{article}
\usepackage[utf8]{inputenc}
\usepackage[T1]{fontenc}
\usepackage{ctex} % 支持中文
\usepackage{geometry}
\usepackage{amsmath, amssymb}
\usepackage{enumitem}
\usepackage{hyperref}
\usepackage{fancyhdr}

% 页面边距设置
\geometry{left=2.5cm, right=2.5cm, top=2.5cm, bottom=2.5cm}

% 页眉页脚设置
\pagestyle{fancy}
\fancyhf{}
\rhead{C语言与数据库原理真题汇编}
\lhead{专升本复习资料}
\rfoot{第 \thepage 页}

\title{\textbf{C语言程序设计专升本真题与数据库原理课后习题全解}}
\author{整理自上传文件}
\date{}

\begin{document}

\maketitle
\tableofcontents
\newpage

% ==============================================================================
% 第一部分：C语言程序设计专升本真题 (基于 c语言程序设计专升本真题附答案.pdf)
% ==============================================================================
\section{C语言程序设计专升本真题 (一)}
\subsection*{一、选择题（共 50 分，每题 2 分）}
\begin{enumerate}
    \item 以下叙述中正确的是
    \begin{enumerate}
        \item[A)] C语言比其他语言高级
        \item[B)] C语言可以不用编译就能被计算机识别执行
        \item[C)] C语言以接近英语国家的自然语言和数学语言作为语言的表达形式
        \item[D)] C语言出现的最晚，具有其他语言的一切优点
    \end{enumerate}
    
    \item C语言中用于结构化程序设计的三种基本结构是
    \begin{enumerate}
        \item[A)] 顺序结构、选择结构、循环结构
        \item[B)] if、switch、break
        \item[C)] for、while、do-while
        \item[D)] if、for、continue
    \end{enumerate}
    
    \item 在一个 C程序中
    \begin{enumerate}
        \item[A)] main函数必须出现在所有函数之前
        \item[B)] main函数可以在任何地方出现
        \item[C)] main函数必须出现在所有函数之后
        \item[D)] main函数必须出现在固定位置
    \end{enumerate}
    
    \item 下列叙述中正确的是
    \begin{enumerate}
        \item[A)] C语言中既有逻辑类型也有集合类型
        \item[B)] C语言中没有逻辑类型但有集合类型
        \item[C)] C语言中有逻辑类型但没有集合类型
        \item[D)] C语言中既没有逻辑类型也没有集合类型
    \end{enumerate}
    
    \item 下列关于 C语言用户标识符的叙述中正确的是
    \begin{enumerate}
        \item[A)] 用户标识符中可以出现在下划线和中划线（减号）
        \item[B)] 用户标识符中不可以出现中划线，但可以出现下划线
        \item[C)] 用户标识符中可以出现下划线，但不可以放在用户标识符的开头
        \item[D)] 用户标识符中可以出现在下划线和数字，它们都可以放在用户标识符的开头
    \end{enumerate}
    
    \item 以下叙述中正确的是
    \begin{enumerate}
        \item[A)] 构成 C程序的基本单位是函数
        \item[B)] 可以在一个函数中定义另一个函数
        \item[C)] main（）函数必须放在其他函数之前
        \item[D)] C函数定义的格式是 K\&R格式
    \end{enumerate}
    
    \item 应用数据库的主要目的是
    \begin{enumerate}
        \item[A)] 解决数据保密问题
        \item[B)] 解决数据完整性问题
        \item[C)] 解决数据共享问题
        \item[D)] 解决数据量大的问题
    \end{enumerate}
    
    \item 一个 C语言程序是由
    \begin{enumerate}
        \item[A)] 一个主程序和若干子程序组成
        \item[B)] 函数组成
        \item[C)] 若干过程组成
        \item[D)] 若干子程序组成
    \end{enumerate}
    
    \item 请选出可用作 C语言用户标识符的是
    \begin{enumerate}
        \item[A)] void,define,WORD
        \item[B)] a3\_b3,\_123,IF
        \item[C)] FOR,--abc,Case
        \item[D)] 2a,Do,Sizeof
    \end{enumerate}
    
    \item 下列各数据类型不属于构造类型的是
    \begin{enumerate}
        \item[A)] 枚举型
        \item[B)] 共用型
        \item[C)] 结构型
        \item[D)] 数组型
    \end{enumerate}
    
    \item 在 16位 C编译系统上，若定义 long a；，则能给 a赋 40000 的正确语句是
    \begin{enumerate}
        \item[A)] a=20000+20000;
        \item[B)] a=4000\*10;
        \item[C)] a=30000+10000;
        \item[D)] a=4000L\*10L
    \end{enumerate}
    
    \item 以下不正确的叙述是
    \begin{enumerate}
        \item[A)] 在 C程序中，逗号运算符的优先级最低
        \item[B)] 在 C程序中，APH和 aph是两个不同的变量
        \item[C)] 若 a和 b类型相同，在计算了赋值表达式 a=b后 b中的值将放入 a中，而 b中的值不变
        \item[D)] 当从键盘输入数据时，对于整型变量只能输入整型数值，对于实型变量只能输入实型数值
    \end{enumerate}
    
    \item sizeof（float）是
    \begin{enumerate}
        \item[A)] 一个双精度型表达式
        \item[B)] 一个整型表达式
        \item[C)] 一种函数调用
        \item[D)] 一个不合法的表达式
    \end{enumerate}
    
    \item 若 x,i,j和 k都是 int型变量，则计算表达式 x= （i=4,j=16,k=32）后，x的值为
    \begin{enumerate}
        \item[A)] 4
        \item[B)] 16
        \item[C)] 32
        \item[D)] 52
    \end{enumerate}
    
    \item 设有定义：int n=0,\*p=\&n,\*\*q=\&p,则下列选项中正确的赋值语句是
    \begin{enumerate}
        \item[A)] p=1；
        \item[B)] \*q=2;
        \item[C)] q=p;
        \item[D)] \*p=5;
    \end{enumerate}
    
    \item 以下叙述中正确的是
    \begin{enumerate}
        \item[A)] C程序的基本组成单位是语句
        \item[B)] C程序中的每一行只能写一条语句
        \item[C)] 简单 C语句必须以分号结束
        \item[D)] C语句必须在一行内写完
    \end{enumerate}
    
    \item 计算机能直接执行的程序是
    \begin{enumerate}
        \item[A)] 源程序
        \item[B)] 目标程序
        \item[C)] 汇编程序
        \item[D)] 可执行程序
    \end{enumerate}
    
    \item 以下关于宏的叙述中正确的是
    \begin{enumerate}
        \item[A)] 宏名必须用大写字母表示
        \item[B)] 宏定义必须位于源程序中所有语句之前
        \item[C)] 宏替换没有数据类型限制
        \item[D)] 宏调用比函数调用耗费时间
    \end{enumerate}
    
    \item 以下选项中正确的定义语句是
    \begin{enumerate}
        \item[A)] double a;b;
        \item[B)] double a=b=7
        \item[C)] double a=7,b=7;
        \item[D)] double,a,b;
    \end{enumerate}
    
    \item 以下不能正确表示代数式 $\frac{2ab}{cd}$ 的 C语言表达式是
    \begin{enumerate}
        \item[A)] 2\*a\*b/c/d
        \item[B)] a\*b/c/d\*2
        \item[C)] a/c/d\*b\*2
        \item[D)] 2\*a\*b/c\*d
    \end{enumerate}
    
    \item C源程序中不能表示的数制是
    \begin{enumerate}
        \item[A)] 二进制
        \item[B)] 八进制
        \item[C)] 十进制
        \item[D)] 十六进制
    \end{enumerate}
    
    \item 检查软件产品是否符合需求定义的过程称为
    \begin{enumerate}
        \item[A)] 确认测试
        \item[B)] 集成测试
        \item[C)] 验证测试
        \item[D)] 验收测试
    \end{enumerate}
    
    \item 数据流图用于抽象描述一个软件的逻辑模型，数据流图由一些特定的图符构成。下列图符名标识的图符不属于数据流图合法图符的是
    \begin{enumerate}
        \item[A)] 控制流
        \item[B)] 加工
        \item[C)] 数据存储
        \item[D)] 源和潭
    \end{enumerate}
    
    \item 若变量已正确定义为 int型，要通过语句 scanf(〞\%d,\%d,\%d〞,\&a,\&b,\&c);给 a赋值 1、给 b赋值 2、给 c赋值 3，以下输入形式中错误的是（u代表一个空格符）
    \begin{enumerate}
        \item[A)] uuu1,2,3<回车>
        \item[B)] 1u2u3<回车>
        \item[C)] 1,uuu2,uuu3<回车>
        \item[D)] 1,2,3<回车>
    \end{enumerate}
    
    \item 有以下程序段
    \begin{verbatim}
int a,b,c;
a=10;b=50;c=30;
if(a>b)a=b,b=c,c=a;
printf(〞a=%d b=%d c=%d\n〞,a,b,c);
    \end{verbatim}
    程序的输出结果是
    \begin{enumerate}
        \item[A)] a=10 b=50 c=10
        \item[B)] a=10 b=50 c=30
        \item[C)] a=10 b=30 c=10
        \item[D)] a=50 b=30 c=50
    \end{enumerate}
\end{enumerate}

\subsection*{二、填空题（共 20 分，每题 2 分）}
\begin{enumerate}
    \item 关系操作的特点是\_\_\_\_\_\_操作。
    \item 按照软件测试的一般步骤，集成测试应在\_\_\_\_\_\_测试之后进行。
    \item 软件工程三要素包括方法、工具和过程，其中，\_\_\_\_\_\_支持软件开发各个环节的控制和管理。
    \item E-mail地址由用户和域名两部分组成，这两部分的分隔符为\_\_\_\_\_\_。
    \item 在二维表中，元组的\_\_\_\_\_\_不能再分成更小的数据项。
    \item 设变量 a和 b已正确定义并赋初值。请写出与 a-=a+b等价的赋值表达式：\_\_\_\_\_\_。
    \item 在 DOS环境下，表示打印机的设备文件名为\_\_\_\_\_\_。
    \item 数据的逻辑结构有线性结构和\_\_\_\_\_\_两大类。
    \item 顺序存储方法是把逻辑上相邻的结点存储在物理位置的\_\_\_\_\_\_存储单元中。
    \item 一个类可以从直接或间接的祖先中继承所有属性和方法。采用这个方法提高了软件的\_\_\_\_\_\_。
\end{enumerate}

\subsection*{三、程序设计填空题（共 30 分，每空 2 分）}
\subsubsection*{1. 输入 a、b、c三个值，输出其中最大者。}
\begin{verbatim}
main()
{int a,b,c,max;
    ("请输入三个数 a,b,b:\n");
    ("%d,%d,%d",&a,&b,&c);
    max=a;
    if(max<b)
        max=;
    if(max < c)
        max=c;
    printf("最大数为：%d",);
}
\end{verbatim}

\subsubsection*{2. 求 1!+2!++20!的值。}
\begin{verbatim}
main()
{s=0,t=1;
    int n;
    for(n=1;n<=20;)
    {
        t=t * n;
        s=s + t;
    }
    ("1!+2!++20!=%e\n",s);
}
\end{verbatim}

\subsubsection*{3. 一球从 100m高度自由落下,每次落地后反跳回原高度的一半,再落下。求它在第 10次落地时,共经过多少 m?第 10次反弹多高?}
\begin{verbatim}
main()
{float sn=100,hn=;
    int n;
    for(n=2;n<=10;n++)
    {sn=sn+;
        hn=;
    }
    printf("第 10次落地时共经过%f m.\n",);
    printf("第 10次反弹%f m.\n",);
}
\end{verbatim}

% ==============================================================================
% 第二部分：C语言程序设计专升本真题 (基于 (完整版)专升本C语言真题.pdf)
% ==============================================================================
\section{C语言程序设计专升本真题 (二)}

\subsection*{05年 C语言}
\subsubsection*{六、单项选择 (10分，每题 1分)}
\begin{enumerate}
    \item 下列变量定义正确的是
    \begin{enumerate}
        \item[A)] int 2ab；
        \item[B)] float a>b；
        \item[C)] char \$123
        \item[D)] int\*per；
    \end{enumerate}
    
    \item 若有说明 int a=1,x=2,y=3；下列不是 C语言合法表达式的是
    \begin{enumerate}
        \item[A)] ++9
        \item[B)] （float)(x)
        \item[C)] a++
        \item[D)] （float)x+y
    \end{enumerate}
    
    \item 若有说明 int i=1,j=2,k=3;表达式 i\&\&j\&\&k的值为
    \begin{enumerate}
        \item[A)] 1
        \item[B)] 2
        \item[C)] 3
        \item[D)] 4
    \end{enumerate}
    
    \item 若有说明 int a,b;下面输入函数调用语句正确的是
    \begin{enumerate}
        \item[A)] scanf（“\%d\%d”,a,b）;
        \item[B)] scanf（“％d\%d”,\&a,\&b)；
        \item[C)] scanf(\%d\%d,a,b）；
        \item[D)] scanf(％d\%d，＆a,\&b);
    \end{enumerate}
    
    \item 下列不是关系表达式的是
    \begin{enumerate}
        \item[A)] 3>5
        \item[B)] 1<2>3
        \item[C)] !3>5
        \item[D)] 1+2>3
    \end{enumerate}
    
    \item 下列对数组的定义正确的是
    \begin{enumerate}
        \item[A)] int a(10)；
        \item[B)] int a［10］;
        \item[C)] int a\{10｝；
        \item[D)] int a10；
    \end{enumerate}
    
    \item 下列对 C程序结构的说法不正确的是
    \begin{enumerate}
        \item[A)] C程序是由一系列函数构成的
        \item[B)] C程序可以有多个 main（）函数
        \item[C)] C程序中函数名不可以和变量名相同
        \item[D)] C程序中可以定义函数
    \end{enumerate}
    
    \item 若有说明 int\*p,a[10],j=3；下列指针变量赋值错误的是
    \begin{enumerate}
        \item[A)] p=\&j；
        \item[B)] p=\&a［j］;
        \item[C)] p=a;
        \item[D)] p=0x1000；
    \end{enumerate}
    
    \item 若有说明 int a=4；执行语句 a>>1后,变量 a的值
    \begin{enumerate}
        \item[A)] 1
        \item[B)] 2
        \item[C)] 3
        \item[D)] 4
    \end{enumerate}
    
    \item 以只读的方式打开文本文件“test.txt”的正确方法是
    \begin{enumerate}
        \item[A)] fopen(“test.txt","r”)；
        \item[B)] fopen（“test.txt","rb”）;
        \item[C)] fopen（“test","r"）;
        \item[D)] fopen（“test.txt")；
    \end{enumerate}
\end{enumerate}

\subsubsection*{七、填空题（10分，每题 1分）}
\begin{enumerate}
    \item a) 若有说明 a=-1；printf（“％d，％x，％o\\n",a,a,a)的输出结果是\_\_\_\_\_\_。
    \item b) -32760在内存中的存储形式是\_\_\_\_\_\_（用十六进制表示)。
    \item c) 7\%4的值为\_\_\_\_\_\_。
    \item d) 写出 C语言中的三种逻辑运算符\_\_\_\_\_\_。
    \item e) 循环语句有 for语句、\_\_\_\_\_\_和\_\_\_\_\_\_。
    \item f) continue语句的作用是\_\_\_\_\_\_。
    \item g) 字符串“123\x45\19abc”的长度为\_\_\_\_\_\_。
    \item h) 数组 a[10］的第 i个元素的指针是\_\_\_\_\_\_。
    \item i) 若有结构体类型定义 struct STU\{int a; float x; char c；｝；sizeof(struct STU)的值是\_\_\_\_\_\_。
    \item j) FILE\*fp;的作用是定义了一个\_\_\_\_\_\_。
\end{enumerate}

\subsubsection*{八、写出下列程序的运行结果（10分，每题 5分）}
\begin{enumerate}
    \item 
    \begin{verbatim}
main()
{ int i,s;
  for(i=10,s=0;i;s+=i,i--)
      printf("result:%d\n",s);
}
    \end{verbatim}
    
    \item 
    \begin{verbatim}
main()
{ void fun();
  float x,y;
  x=10;
  fun(x,&y);
  printf("result:%.0f,%.0f\n",x,y);
}
void fun(x,y)
float x,*y;
{ *y=x*x;
}
    \end{verbatim}
\end{enumerate}

\subsubsection*{九、编写程序（20分）}
\begin{enumerate}
    \item 编写程序求某三位数的平方和（8分）。
    \item 编写程序求数列 1,1,2,3,5,8,13,...的前 100项的和及平均值（12分）。
\end{enumerate}

\subsection*{06年 C语言}
\subsubsection*{六、单项选择 (15分，每题 1分)}
\begin{enumerate}
    \item C语言程序的基本单位是
    \begin{enumerate}
        \item[A)] 程序行
        \item[B)] 语句
        \item[C)] 函数
        \item[D)] 字符
    \end{enumerate}
    
    \item 可用作 C语言用户标识符的一组字符串是
    \begin{enumerate}
        \item[A)] void define WORD
        \item[B)] a3\_b3\_123 IF
        \item[C)] For –abc Case
        \item[D)] 2a DO sizeof
    \end{enumerate}
    
    \item 设 int a=12，则执行完语句 a+=a-=a\*a；后 a的值是
    \begin{enumerate}
        \item[A)] 552
        \item[B)] 264
        \item[C)] 144
        \item[D)] －264
    \end{enumerate}
    
    \item 以下叙述正确的是
    \begin{enumerate}
        \item[A)] do-while语句构成的循环不能用其它语句构成的循环来代替。
        \item[B)] do-while语句构成的循环只能用 break语句退出。
        \item[C)] 用 do-while语句构成的循环，在 while后的表达式为非零时结束循环。
        \item[D)] 用 do-while语句构成的循环，在 while后的表达式为零时结束循环。
    \end{enumerate}
    
    \item 设有说明 int (\*ptr)[10]其中的标识符 ptr是
    \begin{enumerate}
        \item[A)] 10个指向整型变量的指针
        \item[B)] 指向 10个整型变量的函数指针
        \item[C)] 一个指向具有 10个整型元素的一维数组的指针
        \item[D)] 具有 10个指针元素的一维指针数组，每个元素都只能指向整型量
    \end{enumerate}
    
    \item 有以下程序段
    \begin{verbatim}
typedef struct NODE{ int num; struct NODE*next; } OLD;
    \end{verbatim}
    则以下叙述中正确的是
    \begin{enumerate}
        \item[A)] 以上的说明形式非法
        \item[B)] NODE是一个结构体类型
        \item[C)] OLD是一个结构体类型
        \item[D)] OLD是一个结构体变量
    \end{enumerate}
    
    \item 以下不能正确计算代数式 $\frac{\sin^2(0.5)}{3}$ 值的 C语言表达式是
    \begin{enumerate}
        \item[A)] 1/3\*sin(1/2)\*sin(1/2)
        \item[B)] sin(0.5)*sin(0.5)/3
        \item[C)] pow(sin(0.5),2)/3
        \item[D)] 1/3.0*pow(sin(1.0/2),2)
    \end{enumerate}
    
    \item C语言规定，程序中各函数之间
    \begin{enumerate}
        \item[A)] 既允许直接递归调用也允许间接递归调用
        \item[B)] 不允许直接递归调用也不允许间接递归调用
        \item[C)] 允许直接递归调用不允许间接递归调用
        \item[D)] 不允许直接递归调用允许间接递归调用
    \end{enumerate}
    
    \item 在宏定义 \#define PI 3.14159中，用宏名 PI代替一个
    \begin{enumerate}
        \item[A)] 单精度数
        \item[B)] 双精度数
        \item[C)] 常量
        \item[D)] 字符串
    \end{enumerate}
    
    \item 在 C语言中，要求运算数必须是整型的运算符是
    \begin{enumerate}
        \item[A)] \%
        \item[B)] /
        \item[C)] <
        \item[D)] !
    \end{enumerate}
    
    \item 为表示关系 $x \ge y \ge z$，应使用的 C语言表达式是
    \begin{enumerate}
        \item[A)] (x>=y)\&\&(y>=z)
        \item[B)] (x>=y) AND(y>=z)
        \item[C)] (x>=y>=z)
        \item[D)] (x>=y)\&(y>=z)
    \end{enumerate}
    
    \item 有以下程序段
    \begin{verbatim}
int k=0,a=3,b=4,c=5;
k=a>c?c:k;
    \end{verbatim}
    执行该程序段后，k的值是
    \begin{enumerate}
        \item[A)] 3
        \item[B)] 2
        \item[C)] 1
        \item[D)] 0
    \end{enumerate}
    
    \item 若有定义 char\*s="\\\\Name\\\\Address\\n",则指针 s所指字符串长度为
    \begin{enumerate}
        \item[A)] 19
        \item[B)] 15
        \item[C)] 18
        \item[D)] 说明不合法
    \end{enumerate}
    
    \item 下述对 C语言字符数组的描述中错误的是
    \begin{enumerate}
        \item[A)] 字符数组可以存放字符串
        \item[B)] 字符数组中的字符串可以整体输入输出
        \item[C)] 可以在赋值语句中通过赋值运算符对字符数组整体赋值
        \item[D)] 不可以用关系运算符对字符数组中的字符串进行比较
    \end{enumerate}
    
    \item 设有如下的函数 exam(float x){ printf("\%f",x*x); } 则函数的类型为
    \begin{enumerate}
        \item[A)] 与参数 x的类型相同
        \item[B)] 是 void
        \item[C)] 是 int
        \item[D)] 无法确定
    \end{enumerate}
\end{enumerate}

\subsubsection*{七、阅读下列程序，写出其运行结果 (每小题 5分，共 25分)}
\begin{enumerate}
    \item 
    \begin{verbatim}
main(){
    int i,j,x;
    for(i=0;i<=4;i++){
        for(j=1;j<=4-i;j++) printf(" ");
        for(j=0;j<=2*i+1;j++) printf("*");
        printf("\n");
    }
}
    \end{verbatim}
    
    \item 
    \begin{verbatim}
main(){
    int k=3,n=0;
    while(k>0){
        switch(k){
            case 1: n+=k;
            case 2:
            case 3: n+=k;
            default: break;
        }
        k--;
    }
    printf("%d\n",n);
}
    \end{verbatim}
    
    \item 
    \begin{verbatim}
main(){
    int i,j,row,column,m;
    static int array[3][3]={{100,200,300},{28,72,-30},{-850,2,6}};
    m=array[0][0];
    for(i=0;i<3;i++)
        for(j=0;j<3;j++)
            if(array[i][j]<m){
                m=array[i][j];
                row=i; column=j;
            }
    printf("%d,%d,%d\n",m,row,column);
}
    \end{verbatim}
    
    \item 
    \begin{verbatim}
#include<stdio.h>
int p(int k,int a[]){
    int m,i,c=0;
    for(m=2;m<=k;m++)
        for(i=2;i<m;i++){
            if(!(m%i)) break;
            if(i==m) a[c++]=m;
        }
    return c;
}
#define MAXN 20
main(){
    int i,m,s[MAXN];
    m=p(13,s);
    for(i=0;i<m;i++) printf("%4d",s[i]);
    printf("\n");
}
    \end{verbatim}
    
    \item 
    \begin{verbatim}
int f(int n){
    if(n==0||n==1) return 1;
    else return f(n-2)+2*f(n-1);
}
main(){
    int n=5;
    printf("%d",f(n));
}
    \end{verbatim}
\end{enumerate}

\subsubsection*{八、程序填空；按照要求完成下面的程序（函数）（每空 2分，共 10分）}
\begin{enumerate}
    \item 本函数用对分查找法，在以按字母次序从小到大排序的字符数组 list中查找字符 c，若 c在数组中，函数返回字符 c在数组中的下标，否则返回-1。
    \begin{verbatim}
int search(char list[],char c,int len){
    int low,high,k;
    low=0; high=len-1;
    while( (1) ) {
        k=(low+high)/2;
        if( (2) ) return k;
        else if( (3) ) high=k-1;
        else low=k+1;
    }
    return -1;
}
    \end{verbatim}
    
    \item 函数 mycmp(char *s,char *t)的功能是比较字符串 s和 t的大小，当 s等于 t时返回 0，否则返回 s和 t的第一个不同字符的 ASCII码的差值。
    \begin{verbatim}
mycmp(char*s,char*t){
    while(*s==*t){
        if( (4) ) return 0;
        ++s; ++t;
    }
    return( (5) );
}
    \end{verbatim}
\end{enumerate}

\subsection*{07年 C语言}
\subsubsection*{四、填空题（本题20分，每空 2分）}
\begin{enumerate}
    \item C语言中规定，整型常量可以用十进制、八进制和\_\_\_\_\_\_进制形式来表示。
    \item 结构化程序设计中的三种基本结构为：顺序结构、\_\_\_\_\_\_和循环结构。
    \item 在 C语言中，对于负整数，在内存中是以\_\_\_\_\_\_码形式进行存储。
    \item 在 C语言中，若被定义为 int类型的变量，在内存中占用\_\_\_\_\_\_个字节的存储空间。
    \item 已有定义:int a[5]，*p；当执行了 p=\&a[3]；语句时，是将指针变量 p指向了 a数组的第\_\_\_\_\_\_个元素的地址。
    \item 若某变量被定义为 auto变量的存储单元，则将被分配在内存的\_\_\_\_\_\_存储区域。
    \item 在下列给出的字符数组 c，它在内存中所占用的字节数是\_\_\_\_\_\_。
    \begin{verbatim}
char c[]={"C language"};
    \end{verbatim}
    \item 在 C语言中，能够实现循环结构的语句有： while语句、if/goto语句、do-while语句以及\_\_\_\_\_\_语句。
    \item 若有 a=3,b=5;则求 a>b的关系运算结果是\_\_\_\_\_\_。
    \item 若有定义 a[10];则允许数组 a的下标最小可以是\_\_\_\_\_\_。
\end{enumerate}

\subsubsection*{五、请写出下列程序的运行结果（本题 10分，每小题 2分）}
\begin{enumerate}
    \item 
    \begin{verbatim}
main(){
    int n=100;
    if(n>100) printf("***");
    else printf("###");
}
    \end{verbatim}
    
    \item 
    \begin{verbatim}
main(){
    int a=2,b=-1,c=2;
    if(a<b)
        if(b<0) c=0;
        else c+=1;
    printf("c=%d\n",c);
}
    \end{verbatim}
    
    \item 
    \begin{verbatim}
main(){
    char s[]="student\0teacher";
    printf("%s\n",s);
}
    \end{verbatim}
    
    \item 
    \begin{verbatim}
main(){
    int a=3,b=4;
    printf("a=%d,b=%d\n",++a,b++);
}
    \end{verbatim}
    
    \item 
    \begin{verbatim}
main(){
    static int a[5],i;
    for(i=0;i<5;i++) a[i]=a[i]+i;
    for(i=0;i<5;i++) printf("%d,",a[i]);
}
    \end{verbatim}
\end{enumerate}

\subsubsection*{六、单选题（本题 10分，每小题 2分）}
\begin{enumerate}
    \item 
    \begin{verbatim}
main(){
    int k=11;
    printf("k=%d, k=%o, k=%x\n",k,k,k);
}
    \end{verbatim}
    A. k=11, k=12, k=11 \\
    B. k=11, k=13, k=13 \\
    C. k=11, k=013, k=0xb \\
    D. k=11, k=13, k=b
    
    \item 
    \begin{verbatim}
main(){
    int y=10;
    while(y--);
    printf("y=%d\n",y);
}
    \end{verbatim}
    A. y=10 \\
    B. y=1 \\
    C. y=随机值 \\
    D. y=-1
    
    \item 
    \begin{verbatim}
main(){
    int a,b,*p1,*p2;
    p1=&a; p2=&b;
    *p1=100; *p2=200;
    c=*p1+*p2;
    printf("%d\n",c);
}
    \end{verbatim}
    A. 300 \\
    B. 100+200 \\
    C. 100 \\
    D. 200
    
    \item 在下列程序中，当执行到 gets(ss)；语句时，若输入字符为”ABC”时，则该程序的输出结果是：
    \begin{verbatim}
main(){
    char ss[10]="12345";
    strcat(ss,"6789");
    gets(ss);
    printf("%s\n",ss);
}
    \end{verbatim}
    A. ABC \\
    B. ABC9 \\
    C. 123456ABC \\
    D. ABC456789
    
    \item 
    \begin{verbatim}
main(){
    char a[]="morning",t;
    int i,j=0;
    for(i=1;i<7;i++)
        if(a[j]<a[i]) j=i;
    t=a[j]; a[j]=a[7]; a[7]=t;
    puts(a);
}
    \end{verbatim}
    A. mogninr \\
    B. mo \\
    C. morning \\
    D. mornin
\end{enumerate}

\subsubsection*{七、编程题（10分，每题 5分）}
\begin{enumerate}
    \item 请将下列一组数据读入到 S数组中，并从中找出最小的值并输出。
    \begin{verbatim}
30,56,88,45,100,20
    \end{verbatim}
    \item 请将下列给出的字符串读入到 ss数组中，并输出该字符串。
    \begin{verbatim}
Student and Teacher
    \end{verbatim}
\end{enumerate}

\subsection*{08年 C语言}
\subsubsection*{六、填空题（10分，每题 1分）}
\begin{enumerate}
    \item C语言中，基本数据类型包括整型、浮点型和\_\_\_\_\_\_。
    \item 一个 C程序是由\_\_\_\_\_\_组成的。
    \item 在 abc、a\_1、a1b2、auto四个变量中，不合法的是\_\_\_\_\_\_。
    \item 字符串”ab\\\c\n\101"的占用内存的字节数为\_\_\_\_\_\_。
    \item 在运算符+、->、\*=、\&\&中，其优先级最低的是\_\_\_\_\_\_。
    \item 定义共用体类型使用关键字\_\_\_\_\_\_。
    \item C语言中，break语句通常用在\_\_\_\_\_\_语句和循环语句中。
    \item 表达式 7\*2/5+3.5+'b'值的类型是\_\_\_\_\_\_。
    \item 设有语句 int a=5；执行语句 printf("\%d",++a);后，输出结果为\_\_\_\_\_\_。
    \item 执行下面程序段后，输出的结果为\_\_\_\_\_\_。
    \begin{verbatim}
for(i=1;i<5;i++) printf("*");
    \end{verbatim}
\end{enumerate}

\subsubsection*{七、判断题（10分，每题 1分）}
\begin{enumerate}
    \item C语言程序总是从源程序文件中的第一个函数开始执行。( )
    \item 数组名代表数组所占存储区的首地址，其值不可以改变。( )
    \item elseif不属于 C语言关键字（保留字）( )
    \item 指针变量可以加减一个整数。( )
    \item 宏替换不占用运行时间。( )
    \item C语言中转义字符以"\\"开头。( )
    \item C语言规定，函数返回值的类型是由 return语句中的表达式类型决定的。( )
    \item 如果在一个函数中的复合语句中定义了一个变量，则该变量只在该复合语句中有效。( )
    \item C语言中的函数既可以递归定义，又可以嵌套定义。( )
    \item main函数可以有参数。( )
\end{enumerate}

\subsubsection*{八、程序分析题（12分，每题 3分）}
\begin{enumerate}
    \item 指出程序的错误并改正
    \begin{verbatim}
#include <stdio.h>
void main(){
    int a,b,max;
    scanf("%d,%d",&a,&b);
    if(a<b) max=a
    max=b;
    printf("max=%d",max);
}
    \end{verbatim}
    
    \item 写出下面程序的运行结果：
    \begin{verbatim}
#include<stdio.h>
void main(){
    int x[]={0,1,2,3,4,5,6,7,8,9};
    int i,sum=0;
    for(i=0;i<10;i=i+2) sum=sum+x[i];
    printf("%d",sum);
}
    \end{verbatim}
    
    \item 写出程序的运行结果
    \begin{verbatim}
#include<stdio.h>
void main(){
    char s[]="ABC",*p;
    for(p=s;p<s+3;p++) printf("%s\n",p);
}
    \end{verbatim}
    
    \item 写出下面程序的功能
    \begin{verbatim}
void ss(char *s1,char*s2){
    while(*s1!='\0') s1++;
    while(*s2!='\0'){
        *s1=*s2;
        s1++; s2++;
    }
    *s1='\0';
}
    \end{verbatim}
\end{enumerate}

\subsubsection*{九、程序设计题（18分，每题 9分）}
\begin{enumerate}
    \item 从键盘上任意输入一个字符串，统计字符串中大小写英文字母出现的次数。
\end{enumerate}

\subsection*{09年 C语言}
\subsubsection*{六、填空题 (8分，每题 2分)}
\begin{enumerate}
    \item 若 a是 int型变量，且 a=5,则表达式 (a+100)\%2+a/2的值为:\_\_\_\_\_\_。
    \item C语言程序中引用标准输入输出库函数，必须在每个源文件的首部写下 \#include<\_\_\_\_\_\_>。
    \item 若 int型变量占内存 2个字节，double型变量占内存 8个字节，有如下定义: union data { int i; double d; }a; 则变量 a在内存中所占字节数为\_\_\_\_\_\_。
    \item 当文件关闭成功后，fclose函数返回值为\_\_\_\_\_\_。
\end{enumerate}

\subsubsection*{七、阅读程序题（15分，每题 3分）}
\begin{enumerate}
    \item 
    \begin{verbatim}
#include <stdio.h>
main(){
    int i=2,j=3,k;
    k=i+j;
    {
        int k=8;
        if(i=3) printf("%d",k);
        else printf("%d",j);
    }
    printf("%d%d",i,k);
}
    \end{verbatim}
    
    \item 
    \begin{verbatim}
#include <stdio.h>
#define SIZE 8
main(){
    char s[]="GDBFHACE";
    int i,j,t;
    for(i=0;i<SIZE;i++){
        for(j=i+1;j<=SIZE;j++){
            if(s[i]>s[j]){
                t=s[i]; s[i]=s[j]; s[j]=t;
            }
        }
    }
    for(i=0;i<SIZE;i++) printf("%c",s[i]);
}
    \end{verbatim}
    
    \item 
    \begin{verbatim}
#include <stdio.h>
int fun(int a,int b,int*cn,int*dn){
    *cn=a*b+b*b;
    *dn=a*a-b*b;
    a=5; b=6;
}
main(){
    int a=2,b=3,c=5,d=6;
    fun(a,b,&c,&d);
    printf("a=%d,b=%d,c=%d\n",a,b,c,d);
}
    \end{verbatim}
    
    \item 
    \begin{verbatim}
#include <stdio.h>
int fun(int x){
    static y=2;
    y++;
    x+=y;
    return x;
}
main(){
    int k;
    k=fun(3);
    printf("%d,%d\n",k,fun(k));
}
    \end{verbatim}
    
    \item 
    \begin{verbatim}
#include<stdio.h>
main(){
    int s=0,m;
    for(m=7;m>=3;m--)
        switch(m){
            case 1: case 4: case 7:s++;break;
            case 2: case 3: case 6:s+=2;
            case 0: case 5:s+=3;break;
        }
    printf("s=%d\n",s);
}
    \end{verbatim}
\end{enumerate}

\subsubsection*{八、完善程序题（15分，每题 3分）}
\begin{enumerate}
    \item 下面程序的功能是找出 100到 200之间不能被 3整除但能被 5整除的数。
    \begin{verbatim}
#include <stdio.h>
main(){
    int m;
    for(m=100;m<=200;m++)
        if (___________);
        printf("%d\t",m);
}
    \end{verbatim}
    
    \item 下面程序通过指向整型变量的指针将数组 m[4][3]的内容按 4行 3列的格式输出，请输出 printf()填入适当的参数，使之通过指针 p将数组元素按要求输出。
    \begin{verbatim}
#include <stdio.h>
main(){
    int m[4][3]={{1,2,3},{4,5,6},{7,8,9},{10,11,12}};
    int i,j,*p=m;
    for(i=0;i<4;i++){
        for(j=0;j<3;j++) printf("%4d",____________);
        printf("\n");
    }
}
    \end{verbatim}
    
    \item 下面程序能够完成交换数组 a和数组 b中的对应元素的功能。
    \begin{verbatim}
#include<stdio.h>
swap(int*p1,int *p2){
    int temp;
    ___________;
}
main(){
    int a[5]={1,3,5,7,9};
    int b[5]={2,4,6,8,10};
    int i;
    for(i=0;i<5;i++) swap(&a[i],&b[i]);
    for(i=0;i<5;i++) printf("a[%d]=%-4d",i,a[i]);
    printf("\n");
    for(i=0;i<5;i++) printf("b[%d]=%-4d",i,b[i]);
    printf("\n");
}
    \end{verbatim}
    
    \item 在某大学举行的演讲比赛中，有十个评委为参赛的选手打分，分数为0～100分。选手最后得分为：去掉一个最高分和一个最低分后其余八个分数的平均值。
    \begin{verbatim}
#include <stdio.h>
main(){
    int max,min,score,i;
    int sum=0;
    max=0; min=100;
    for(i=0;i<10;i++){
        printf("Input number %d=",i++);
        scanf("%d",&score);
        sum+=score;
        if(max<score) max=score;
        if(min>score) min=score;
        printf("\n");
    }
    printf("Canceled max score:%d\nCanceled min score:%d",max,min);
    printf("average score:%lf\n",________________ );
}
    \end{verbatim}
    
    \item 把指针 str所指的字符串按相反的顺序赋给 rev[]。
    \begin{verbatim}
#include<stdio.h>
main(){
    char *str="abcdefg";
    char rev[10];
    int i;
    printf("\n");
    for(i=0;i<7;i++) ______________;
    rev[i]='\0';
    printf("%s\n",rev);
}
    \end{verbatim}
\end{enumerate}

\subsubsection*{九、编程改错题 (12分，每题 3分)}
\begin{enumerate}
    \item 
    \begin{verbatim}
(1) #include<stdio.h>
(2) char a="Beijing";
(3) main( )
(4) {
(5)  printf("%s is one city in china.\n",a);
(6)  p1();
(7)  p2( );
(8) }
(9) p1()
(10){
(11)  printf("%s is one of the biggest city",a);
(12)  return;
(13)}
(14) p2( )
(15){
(16)  printf("in the world.\n");
(17)  return;
(18)}
    \end{verbatim}
    错误的行是：\_\_\_\_\_\_ 改为：\_\_\_\_\_\_
    
    \item 求 $\sum_{k=1}^{100} \sum_{k=1}^{50} \sum_{k=1}^{10} \frac{k+k+k}{k+k+k}$ (注：原题公式显示可能有误，此处按代码逻辑修正)
    \begin{verbatim}
(1) #include <stdio.h>
(2) main()
(3){
(4) int n1=100,n2=50,n3=10;
(5) int k;
(6) float s1=0,s2=0,s3=0;
(7) for(k=1;k<=n1;k++);
(8)  s1=s1+k;
(9) for(k=1;k<=n2;k++);
(10) s2=s2+k*k;
(11) for(k=1;k<=n3;k++);
(12) s3=s3+k/10;
(13) printf("total=%8.2f\n",s1+s2+s3);
(14)}
    \end{verbatim}
    错误的行是：\_\_\_\_\_\_ 改为：\_\_\_\_\_\_
    
    \item 本程序能够在屏幕中央显示出如下图形。
    \begin{verbatim}
#######
#####
###
#
    \end{verbatim}
    \begin{verbatim}
(1) #include <stdio.h>
(2) void main()
(3){
(4)  int i,j,k;
(5)  for(i=1;i<=4;i++)
(6) {
(7)  for(k=1;k<=i+1;k++)
(8) printf(" ");
(9) for(j=1;j<=i;j++)
(10) printf("#");
(11) printf("\n");
(12) }
(13)}
    \end{verbatim}
    错误的行是：\_\_\_\_\_\_ 改为：\_\_\_\_\_\_
    
    \item 将给定的字符串写入指定的文件中去。
    \begin{verbatim}
(1) #include <stdio.h>
(2) main()
(3){
(4)  int i;
(5)  FILE*fp;
(6)  char a[][8]={"Turbo C", "BASIC", "FORTRAN", "COBOL"};
(7)  if(fp=fopen("myfile.txt", "w")==NULL)
(8) {
(9)  printf("can't creat the file.\n");
(10) exit(1);
(11) }
(12) for(i=0;i<4;i++)
(13) {
(14) fputs(a[i],fp);
(15) fputs("\n",fp);
(16) }
(17) fclose(fp);
(18)}
    \end{verbatim}
    错误的行是：\_\_\_\_\_\_ 改为：\_\_\_\_\_\_
\end{enumerate}

\end{document}